Disposem d'un tub de buit com el de la figura. L'elèctrode A és fet de potassi, que té $W_0 = 2,29\ \mathrm{eV}$ com a valor de treball d'extracció.

\begin{apartat}{1}
Determineu la velocitat amb què surten els electrons arrancats de l'elèctrode A quan l'il·luminem amb llum de color violat de 400~nm de longitud d'ona.
\end{apartat}

\figura[0.45]{fig1.png}

\begin{apartat}{1}
A continuació canviem l'elèctrode A per un altre que és fet d'un material desconegut. Per tal de determinar de quin material es tracta, l'il·luminem un altre cop amb la mateixa llum d'abans, i determinem que el potencial de frenada dels electrons de l'elèctrode A és $V_\mathrm{f} = 0,17\ \mathrm{V}$. Determineu el treball d'extracció del material i indiqueu de quin element és fet a partir de la taula de valors següent:

\begin{center}
\begin{tabular}{|c|c|c|c|c|c|c|c|c|c|}
\hline
\textit{Element} & Ba & Li & Mg & As & Al & Bi & Cr & Ag & Be\\
\hline
$W_0$(eV) & 2,70 & 2,93 & 3,66 & 3,75 & 4,08 & 4,34 & 4,50 & 4,73 & 4,98\\
\hline
\end{tabular}
\end{center}
\end{apartat}

\dades{%
  Massa de l'electró, $m_{\text{electró}} = 9,11\times 10^{-31}\ \mathrm{kg}$\\
  Constant de Planck, $h = 6,63\times 10^{-34}\ \mathrm{J\,s}$\\
  Velocitat de la llum, $c = 3,00\times 10^{8}\ \mathrm{m\,s^{-1}}$\\
  $1\ \mathrm{eV} = 1,60\times 10^{-19}\ \mathrm{J}$}
